\section{Integral stability in a fixed number of coordinates}
\label{sec:stability}

The detection argument used additional coordinates in the stable
Steinberg group. To apply Theorem~\ref{thm:detection} to a finitely
presented group, we must now show that adjoining these coordinates
does not kill an element already in a fixed rank. We prove that rank
twelve suffices. All homology in this section has integral coefficients.

Throughout this Section, let $B=\F_2\oplus I_B$ be a unital associative
algebra, with $I_B$ a two-sided ideal such that every finite square
matrix over $I_B$ is nilpotent. No grading or finite generation
assumption is needed. We often suppress $B$ from the notation for its
matrix and Steinberg groups.

\begin{proposition}[Finite-rank inputs]\label{prop:finite-rank-input}
For such an algebra $B$, the following statements hold.
\begin{enumerate}[label=\textup{(\roman*)}]
\item Every tuple of columns of $B^m$ with linearly independent
reductions in $\F_2^m$ extends to an invertible matrix.
\item The kernel $J_r(B)$ is central in $\St_r(B)$ for $r\geq6$.
\item For $r\geq5$, the group $\E_r(B)$ is normal in $\GL_r(B)$,
and the natural maps
\[
 \GL_r(B)/\E_r(B)\longrightarrow
 \GL_{r+1}(B)/\E_{r+1}(B)
\]
are isomorphisms.
\end{enumerate}
\end{proposition}

These are, respectively, Lemma~5.1, Proposition~5.4, and Lemma~5.5
of the companion paper~\cite{Companion}. They have exactly the
hypotheses stated above. The integral conclusion we need is the
following.

\begin{theorem}\label{thm:stability}
Let $B=\F_2\oplus I_B$ satisfy the hypotheses of
Proposition~\ref{prop:finite-rank-input}. For every $n\geq12$, the
standard stabilization map $J_n(B)\to J_{n+1}(B)$ is an isomorphism,
and the natural homomorphism
\[
 \St_n(B)\longrightarrow\St(B)
\]
is injective.
\end{theorem}

The ring $B$ is local: its elements outside $I_B$ are units, whereas
its elements in $I_B$ are nilpotent. Thus the injectivity conclusion
also follows, with a better rank bound, from the classical stability
theorem of Suslin and Tulenbaev~\cite[Theorem~4.1]{SuslinTulenbaev1976}.
We include a direct proof to make the integral homology comparison
and its natural maps explicit. It compares successive ranks by the
action on ordered frames, following the classical method of
homological stability for linear groups; see
van der Kallen~\cite{vanderKallen1980}.
Two details matter here: we use ordered frames with integral chains,
and identify the second homology of their stabilizers by universal
central extensions. We give both arguments explicitly.

\subsection{Central extensions of frame stabilizers}

A group is \emph{perfect} if it equals its commutator subgroup.
A \emph{universal central extension} of a perfect group $H$ is a central
extension $\widetilde H\to H$ admitting a unique homomorphism over $H$
to every central extension of $H$. Its kernel is naturally
$H_2(H;\Z)$; this identification is natural for homomorphisms lifted
between universal central extensions. We use this standard connection
between Steinberg groups and second homology in its integral form;
see~\cite[Theorems~III.5.4 and III.5.5]{WeibelK}.

For $r\geq6$ and $l\geq0$, put
\begin{equation}\label{eq:frame-stabilizer}
 S_{r,l}=
 \left\{\begin{pmatrix}A&0\\ C&\Id_l\end{pmatrix}:
 A\in\E_r(B),\ C\in\Mat_{l\times r}(B)\right\}.
\end{equation}
This is the semidirect product of its additive lower-left strip with
$\E_r(B)$. Replace $\E_r(B)$ in this semidirect product by
$\St_r(B)$, acting through its matrix image, and denote the resulting
group by $\widetilde S_{r,l}$. There is a natural surjection
$\widetilde S_{r,l}\to S_{r,l}$ with kernel $J_r(B)$.

\begin{lemma}\label{lem:stabilizer-uce}
The map $\widetilde S_{r,l}\to S_{r,l}$ is a universal central
extension. Consequently the inclusion of the complementary block
induces a natural isomorphism
\begin{equation}\label{eq:stabilizer-homology}
 H_2(\E_r(B);\Z)\xrightarrow{\ \sim\ }
 H_2(S_{r,l};\Z),
\end{equation}
and both groups identify naturally with $J_r(B)$.
\end{lemma}

\begin{proof}
The extension is central by
Proposition~\ref{prop:finite-rank-input}\textup{(ii)}, because its
kernel acts trivially on the strip. Both its source and target are
perfect. Indeed, an internal root is a consecutive-root commutator,
and a strip root has the same expression using an intermediate index
among the first $r$ coordinates.

Consider an arbitrary central extension $T\to S_{r,l}$. The roots
available in $S_{r,l}$ are $e_{ij}(a)$ with $j\leq r$ and
$i\ne j$, where $i$ may also be one of the last $l$ coordinates.
Choose arbitrary lifts $y_{ij}(a)\in T$. We first show that lifts of
roots appearing in a Steinberg commuting relation actually commute:
\begin{equation}\label{eq:lifts-commute}
 [y_{ij}(a),y_{kt}(b)]=1
 \qquad\text{if }j\ne k\text{ and }i\ne t.
\end{equation}
Choose $h\leq r$ distinct from $i,j,k,t$. In matrices,
\[
 e_{ij}(a)=[e_{ih}(a),e_{hj}(1)],
\]
and both factors on the right commute with $e_{kt}(b)$. In $T$,
their commutators with $y_{kt}(b)$ are therefore central. Conjugating
each factor by $y_{kt}(b)$ changes it only by a central factor, so
fixes their commutator. Since this commutator differs from
$y_{ij}(a)$ by a central element, \eqref{eq:lifts-commute} follows.
The bound $r\geq6$ supplies every spare index used below as well.

Define corrected lifts by
\begin{equation}\label{eq:canonical-root-lift}
 \widehat y_{ij}(a)=[y_{ih}(a),y_{hj}(1)],
 \qquad h\leq r,\quad h\ne i,j.
\end{equation}
They are independent of $h$ and satisfy
\begin{equation}\label{eq:lifts-product}
 [y_{ih}(a),y_{hj}(b)]=\widehat y_{ij}(ab)
 \qquad(h\leq r;\ i,h,j\text{ distinct}).
\end{equation}
To prove this, choose a fourth distinct index $f\leq r$ and apply
Lemma~\ref{lem:commutator} to
\[
 X=y_{ih}(a),\qquad Y=y_{hf}(b),\qquad Z=y_{fj}(1).
\]
All its commuting hypotheses follow from
\eqref{eq:lifts-commute}, since $[X,Y]$ and $[Y,Z]$ lift,
respectively, $e_{if}(ab)$ and $e_{hj}(b)$. The resulting identity
equates the left side of \eqref{eq:lifts-product} with
$[y_{if}(ab),y_{fj}(1)]$: replacing a lift by a central multiple does
not change a commutator. Setting $b=1$ proves independence in
\eqref{eq:canonical-root-lift}; the same identity then proves
\eqref{eq:lifts-product} for arbitrary $b$.

The corrected lifts are additive in their parameters. In
\eqref{eq:canonical-root-lift}, use
$y_{ih}(a)y_{ih}(a')$ as a lift of $e_{ih}(a+a')$ and expand its
commutator with $y_{hj}(1)$. Equation~\eqref{eq:lifts-commute}
allows the commutator factors to pass the $ih$ factors, giving
\[
 \widehat y_{ij}(a+a')=
 \widehat y_{ij}(a)\widehat y_{ij}(a').
\]
Here $\widehat y_{ij}(0)=1$, since any lift of an identity matrix
is central. Together with \eqref{eq:lifts-commute} and
\eqref{eq:lifts-product}, this proves all Steinberg relations for
the corrected lifts on the available roots.

On the first $r$ coordinates they therefore define a homomorphism
$\St_r(B)\to T$. On the strip they define a homomorphism from the
additive group $\Mat_{l\times r}(B)$ to $T$. These homomorphisms
respect the semidirect-product action: conjugation of a strip root
by an internal root is determined by a commuting relation, except
when their indices form a consecutive path, in which case it is
determined by \eqref{eq:lifts-product}. These are exactly the
matrix conjugation rules. Thus they give a homomorphism
$\widetilde S_{r,l}\to T$ over $S_{r,l}$. It is unique, since two
such homomorphisms differ by a homomorphism from the perfect group
$\widetilde S_{r,l}$ to the central, abelian kernel of $T\to S_{r,l}$.
This proves universality.

Taking $l=0$ identifies $H_2(\E_r(B);\Z)$ with $J_r(B)$.
The inclusion $\E_r(B)\to S_{r,l}$ lifts by adjoining the zero
strip and is the identity on the kernels $J_r(B)$. Naturality of
the universal central extension now proves
\eqref{eq:stabilizer-homology}.
\end{proof}

\subsection{The integral complex of ordered frames}

An \emph{ordered frame} of $B^m$ is an ordered tuple of columns
whose reductions modulo $I_B$ are linearly independent. Let $C_p(m)$
be the free abelian group on ordered frames of length $p+1$, for
$p\geq0$, with boundary the alternating sum of deletions. Set
$C_{-1}(m)=\Z$, generated by the empty frame, and augment by sending
each one-column frame to $1$. Every ordering is a separate generator;
we impose no relation that identifies two orderings by a sign.

The corresponding finite complexes are complexes of injective words
on matroid independence complexes. Their shellability is established
in~\cite[Theorem~1.2 and Example~3.3]{JonssonWelker2009}. The following
direct proof gives the precise integral vanishing range we need.

\begin{lemma}\label{lem:ordered-frames}
The augmented chain complex $C_\bullet(m)$ is exact in degrees at
most $m-2$.
\end{lemma}

\begin{proof}
First let $V$ be a finite labeled collection of nonzero vectors
spanning a vector space of dimension $m$. Different labels remain
different vertices even if their vectors agree. Write $C_\bullet(V)$
for the augmented complex on its ordered independent tuples. We prove
the claimed vanishing by induction on the number of labels, allowing
$m=0$: the empty collection has its only chain group in degree $-1$,
and the asserted range is then empty.

Delete a label $v$, and let $V'=V\setminus\{v\}$. In the relative
complex
\[
 Q_\bullet=C_\bullet(V)/C_\bullet(V'),
\]
the generators are the independent tuples containing $v$.
Filter $Q_\bullet$ increasingly by the number $a$ of labels before
$v$. Deleting $v$ is zero in the relative complex; deleting a label
before it lowers the filtration; deleting a label after it preserves
the filtration. Thus the associated graded complex for a fixed $a$
is a direct sum, over independent ordered prefixes
$(u_1,\ldots,u_a,v)$, of complexes in which only the suffix varies.

For one such prefix, project all remaining labeled vectors to the
quotient by $\langle u_1,\ldots,u_a,v\rangle$, and discard the zero
projections. The projected collection spans dimension $m-a-1$ and
has fewer labels than $V$. Its augmented ordered-tuple complex,
shifted up by $a+1$ degrees, is precisely the corresponding summand
of the associated graded complex, up to the constant boundary sign
$(-1)^{a+1}$. By induction its homology vanishes through degree
\[
 (m-a-1)-2+(a+1)=m-2.
\]
This also covers a zero-dimensional quotient: its sole homology in
degree $-1$ shifts to degree $a=m-1$, outside the asserted range.
The filtration is finite, so $H_i(Q_\bullet)=0$ for $i\leq m-2$.

If $V'$ still spans dimension $m$, induction and the long exact
sequence of the pair give the required vanishing for $V$.
Otherwise $V'$ spans dimension $m-1$, and $v$ is independent of its
entire span. Prepending $v$ to every tuple of $V'$ gives a chain
homotopy from its inclusion into $C_\bullet(V)$ to zero; in
augmented degree $-1$ it sends $1$ to $(v)$. The induced maps
$H_i(C_\bullet(V'))\to H_i(C_\bullet(V))$ are therefore zero.
They are also surjective for $i\leq m-2$, by the relative vanishing
just proved. The homology of $C_\bullet(V)$ vanishes in this range.

Finally, any cycle in $C_\bullet(m)$ uses finitely many columns.
Adjoin the standard basis columns to this finite collection. Their
reductions, with column labels retained, form a finite collection
spanning $\F_2^m$. The finite result supplies a bounding chain in
the asserted range, proving the Lemma.
\end{proof}

\subsection{Second homology in successive ranks}

\begin{proof}[Proof of Theorem~\ref{thm:stability}]
We now combine the stabilizer calculation with the ordered-frame
complex to prove Theorem~\ref{thm:stability}. Fix $m\geq13$ and let
$\E_m(B)$ act on ordered frames by left multiplication.
For frame lengths $l\leq5$, this action is transitive. To see this,
complete a given frame to an invertible matrix, placing it in the
last $l$ columns, and put $r=m-l$. By
Proposition~\ref{prop:finite-rank-input}\textup{(iii)}, the map
$\GL_r/\E_r\to\GL_m/\E_m$ is an isomorphism. Right multiplication
by a suitable complementary $\GL_r$ block therefore brings the
completion into $\E_m$ without changing its last $l$ columns.

The pointwise stabilizer of the last $l$ standard columns is exactly
$S_{m-l,l}$ from \eqref{eq:frame-stabilizer}. Indeed, its general
form has $A\in\GL_{m-l}$ and arbitrary lower-left strip. The strip
is elementary, and the injectivity of
$\GL_{m-l}/\E_{m-l}\to\GL_m/\E_m$ says that this matrix belongs
to $\E_m$ precisely when $A\in\E_{m-l}$. Write $S_l=S_{m-l,l}$
for the rest of the proof. All these complementary ranks are at
least eight, so Lemma~\ref{lem:stabilizer-uce} applies.

Take a free right $\Z[\E_m]$-resolution $P_\bullet\to\Z$ of the
trivial module and form the first-quadrant double complex
\begin{equation}\label{eq:frames-double-complex}
 D_{p,q}=P_q\otimes_{\Z[\E_m]}C_p(m),\qquad p,q\geq0.
\end{equation}
This is the usual equivariant homology double complex;
see~\cite[Chapter~VII]{Brown}.
Its total homology in degree two is $H_2(\E_m;\Z)$ via augmentation.
For completeness, adjoining the augmentation column $p=-1$ gives
the same tensor product with the augmented frame complex. Each
$P_q$ is free and hence flat, so
Lemma~\ref{lem:ordered-frames} makes the rows exact in the stated
range. The augmented total complex has zero homology in degrees
one and two. The long exact sequence separating its augmentation
column from \eqref{eq:frames-double-complex} consequently gives
the asserted isomorphism.

Take vertical homology first. For $0\leq p\leq4$, the permutation
module on ordered $(p+1)$-frames has one orbit and stabilizer
$S_{p+1}$. Shapiro's Lemma identifies the first page as
\begin{equation}\label{eq:frames-first-page}
 E^1_{p,q}=H_q(S_{p+1};\Z),\qquad 0\leq p\leq4.
\end{equation}
The first differential is the alternating sum of the face maps,
each expressed by a stabilizer inclusion and a coordinate change
that returns the remaining frame to the chosen standard frame.

The row $q=1$ is zero in these columns, because every $S_l$ is
perfect. In row $q=0$, all entries are $\Z$ and each face map is
the identity. Hence the first differential is zero for odd $p$ and
the identity for even $p$. In particular,
\begin{equation}\label{eq:vanishing-second-page}
 E^2_{1,1}=E^2_{2,1}=E^2_{2,0}=E^2_{3,0}=0.
\end{equation}

In row $q=2$, the two face maps from column one to column zero
coincide. Indeed, the inclusion $\E_{m-2}\to S_2$ induces an
isomorphism on $H_2$ by Lemma~\ref{lem:stabilizer-uce}, so it
suffices to compare the faces on this complementary block.
The first standard frame is $(e_{m-1},e_m)$. Deleting its first
column gives $e_m$. Deleting its second column gives $e_{m-1}$,
which is returned to $e_m$ by swapping the last two coordinates.
This swap is elementary in characteristic two and centralizes the
first $m-2$ coordinates. The two induced maps on that block are
therefore identical. Their alternating sum vanishes, and
\[
 E^2_{0,2}=H_2(S_1;\Z).
\]

No higher differential leaves position $(0,2)$. The only possible
incoming differentials are from $(2,1)$ on the second page and
$(3,0)$ on the third page; both vanish by
\eqref{eq:vanishing-second-page}. The other positions of total
degree two, $(1,1)$ and $(2,0)$, also vanish. The spectral sequence
thus shows that the edge map
\[
 H_2(S_1;\Z)\longrightarrow H_2(\E_m;\Z)
\]
is an isomorphism. This is the map induced by inclusion: the
zero-th column corresponds to the standard column $e_m$, and its
map into the total complex followed by augmentation is the usual
homology map for its stabilizer inclusion.

Combining with \eqref{eq:stabilizer-homology}, we have proved that
the natural stabilization map
\[
 H_2(\E_{m-1}(B);\Z)\longrightarrow H_2(\E_m(B);\Z)
\]
is an isomorphism for $m\geq13$. The natural identifications by
universal central extensions turn this into the isomorphism
$J_{m-1}(B)\to J_m(B)$.

Finally, an element killed by $\St_n(B)\to\St(B)$ already has
identity matrix image, since coordinate embeddings of matrix groups
are injective. It therefore lies in $J_n(B)$. If $n\geq12$, the
successive isomorphisms on these kernels show that it cannot become
trivial at any later rank. The stable group is the direct limit, so
this proves the injectivity asserted in
Theorem~\ref{thm:stability}.
\end{proof}
